\section{M\'ethode}
La m\'ethode glaciologique combine les relev\'es ponctuels et une
interpolation par bandes d'altitude. Chaque bande re\c{c}oit le
bilan moyen de ses balises, pond\'er\'e par sa surface.

\subsection{Interpolation}
Pour une bande $i$ de surface $S_i$ et de bilan $b_i$, le bilan
sp\'ecifique du glacier vaut
\[
  B = \frac{\sum_i S_i\, b_i}{\sum_i S_i}.
\]
Cette moyenne pond\'er\'ee suppose que le bilan varie peu \`a
l'int\'erieur d'une bande, hypoth\`ese raisonnable pour des bandes
de cent m\`etres.

\begin{table}[h]
  \centering
  \begin{tabular}{lrr}
    Bande & Surface (km$^2$) & Bilan (m w.e.) \\
    2800--2900 & 0.42 & $-2.1$ \\
    2900--3000 & 0.77 & $-1.4$ \\
    3000--3100 & 0.91 & $-0.6$ \\
  \end{tabular}
  \caption{Bilan par bande d'altitude.}
\end{table}

\subsection{Incertitudes}
Les erreurs de mesure des balises sont de l'ordre de dix
centim\`etres d'\'equivalent en eau\footnote{Valeur typique pour
des balises en bois relev\'ees au ruban.}. L'erreur
d'interpolation domine largement pour les bandes sans balise.

\newpage
\subsection{Validation}
Une validation g\'eod\'esique, par diff\'erence de mod\`eles
num\'eriques de terrain, compl\`ete la m\'ethode glaciologique
sur des p\'eriodes de cinq \`a dix ans.

\begin{figure}[h]
  \centering
  \rule{30mm}{12mm}\hspace{4mm}%
  \fbox{\parbox{36mm}{Balise B7, relev\'ee
    en septembre, fond\'ee dans la glace.}}
  \caption{Sch\'ema d'une balise.}
\end{figure}

\noindent\begin{minipage}[t]{0.45\linewidth}
  Colonne de gauche : les mesures
  de l'ann\'ee hydrologique.
\end{minipage}\hfill
\begin{minipage}[t]{0.45\linewidth}
  Colonne de droite : la comparaison
  avec la s\'erie longue.
\end{minipage}
